\section{Technical details and proofs}\label{app:technical-details}\label{app:convergence-proofs}
\subsection{Loss estimates}\label{app:proofs}
\begin{lemma}[Smooth PL consequences]\label{lem:pl}
Under Assumption~\ref{ass:loss}, the minimizer set
$\cS_\star:=\{S:\ell(S)=\ell_\star\}$ is nonempty, and
\begin{equation}\label{eq:pl-consequences}
\norm{\nabla\ell(S)}_{\F}^2\le2L(\ell(S)-\ell_\star),\qquad
\dist(S,\cS_\star)^2\le2(\ell(S)-\ell_\star)/\kappa.
\end{equation}
For every perturbation $H$,
\begin{equation}\label{eq:perturb}
\sqrt{\ell(S+H)-\ell_\star}
\le\sqrt{\ell(S)-\ell_\star}+\sqrt{L/2}\norm H_{\F}.
\end{equation}
\end{lemma}

\begin{proof}[Proof of Lemma~\ref{lem:pl}]

Apply the descent lemma at $S-\nabla\ell(S)/L$ for the first inequality.
Since $\nabla\ell$ is globally Lipschitz, the Euclidean gradient flow
$\dot S=-\nabla\ell(S)$ exists for all $t\ge0$.
Writing $h:=\ell(S)-\ell_\star$ gives $\dot h=-\norm{\nabla\ell}^2$ and,
where $h>0$, $-\frac{\dd}{\dd t}\sqrt h\ge\sqrt{\kappa/2}\norm{\dot S}$.
Thus its total remaining length is at most $\sqrt{2h/\kappa}$.
The flow has a limit, with loss gap zero by PL, proving attainment and the distance
bound. Finally, the descent lemma and the first inequality bound the
perturbed loss gap by $(\sqrt h+\sqrt{L/2}\norm H_{\F})^2$.
\end{proof}
\begin{remark}[Quadratic growth and convexity]
Lemma~\ref{lem:pl} proves the quadratic-growth bound
\begin{equation}\label{eq:quadratic-growth}
 \ell(S)-\ell_\star\ge\frac\lambda2\dist(S,\cS_\star)^2
 \qquad\text{with }\lambda=\kappa.
\end{equation}
If a larger quadratic-growth constant $\lambda$ is known, it improves
the distance bounds while leaving the loss rates, which use $\kappa$, unchanged.
The PL-to-quadratic-growth implication does not need convexity
\citep{KarimiNutiniSchmidt2016}. Conversely, convex quadratic growth
with constant \(\lambda\) implies PL with constant \(\lambda/4\):
for a nearest minimizer \(S_\star\),
\(\Delta\le\ip{\nabla\ell(S)}{S-S_\star}
\le\norm{\nabla\ell(S)}_{\F}\sqrt{2\Delta/\lambda}\).
The best PL and quadratic-growth constants need not coincide.
For a nonzero quadratic Hessian $\quadH\succeq0$,
$\lambda=\kappa=\lambda_{\min}^{+}(\quadH)$ is valid and
$\cS_\star=\Svu_\star+\kerop\quadH$.
\end{remark}

\subsection{Covariance closure and conservation laws}
\label{app:covariance-laws}
Direct matrix differentiation proves closure and conservation at
every width. The measure interpretation is given separately in
Appendix~\ref{app:mean-field}.

\begin{proof}[Proof of Proposition~\ref{prop:covariance-closure-laws}]
Since $\dot W_t=-G_tW_t$ and $G_t^\top=G_t$,
differentiating $W_tW_t^\top/n$ gives \eqref{eq:bures-G}.
For any covariance solution, define
$\dot\Phi_t=-G_t\Phi_t$, $\Phi_0=\Id$. Uniqueness for the resulting
linear covariance equation yields
$\Sigma_t=\Phi_t\Sigma_0\Phi_t^\top$, so positive semidefiniteness
is preserved. The identity $G_tJ+JG_t=0$ gives
$\frac{\dd}{\dd t}(\Phi_t^\top J\Phi_t)=0$.
Invertibility of $\Phi_t$ then gives both identities in
\eqref{eq:orbit}. Equivalently,
$\frac{\dd}{\dd t}(J\Sigma_t)=[G_t,J\Sigma_t]$, so
$\frac{\dd}{\dd t}\tr((J\Sigma_t)^k)=0$ for every $k\ge1$.
The first $m$ trace powers determine the characteristic polynomial.
For later use, the block equations are
\begin{equation}\label{eq:block-flow-general}
 \begin{aligned}
 \dSuu_t&=-\Suv_tE_t-E_t^\top\Svu_t,\quad
 \dSvv_t=-\Svu_tE_t^\top-E_t\Suv_t,\quad
 \dSvu_t&=-E_t\Suu_t-\Svv_tE_t.
 \end{aligned}
\end{equation}
\end{proof}

\begin{proof}[Proof of Proposition~\ref{lem:margins}]
For any $c\in\R$, the characteristic congruence gives
\[
 \Sigma_t-cJ=\Phi_t(\Sigma_0-cJ)\Phi_t^\top.
\]
Invertibility of $\Phi_t$ preserves the entire feasible interval of $c$,
so both endpoints are conserved even when $\Sigma_0$ is singular.
Taking principal blocks proves the hidden bounds.
\end{proof}

\begin{proof}[Spectral formula \eqref{eq:gap-spectrum}]
The feasible sets in Definition~\ref{def:support-gap} contain zero, are
closed intervals, and are bounded above by the corresponding hidden-block
minimum eigenvalues. Thus the maxima exist.
For $\Sigma\succ0$, put $H:=\Sigma^{1/2}J\Sigma^{1/2}$.
Sylvester's inertia theorem gives $d_v$ negative and $d_u$ positive
eigenvalues. Since $H^{-1}=\Sigma^{-1/2}J\Sigma^{-1/2}$,
$\Sigma-cJ=\Sigma^{1/2}(\Id-cH^{-1})\Sigma^{1/2}$ gives
$\aplus{\Sigma}=\lambda_{d_v+1}(H)$ and
$\aminus{\Sigma}=-\lambda_{d_v}(H)$.
For singular $\Sigma$, apply this to $\Sigma+\epsilon\Id$.
As $\epsilon\downarrow0$, each margin decreases to a finite limit that
is feasible for $\Sigma$ by closedness, and hence equals its margin.
Continuity of the square root and of ordered symmetric eigenvalues
proves \eqref{eq:gap-spectrum} without deleting any zero eigenvalues.
Finally $H$ and $J\Sigma$ have the same characteristic polynomial by
the $AB,BA$ identity, even when they are not similar.
\end{proof}

\begin{remark}[Sharpness of the supporting constants]
For every $\Sigma\succeq0$, testing the supporting inequalities on the two
coordinate subspaces gives the upper bounds in
\[
 \lambda_{\min}(\Sigma)\le\aplus{\Sigma}\le\lambda_{\min}(\Suu),\qquad
 \lambda_{\min}(\Sigma)\le\aminus{\Sigma}\le\lambda_{\min}(\Svv).
\]
The lower bounds follow from $\Id\mp J\succeq0$.
When $\Svu=0$, the two upper bounds are attained.
Consequently $2\lambda_{\min}(\Sigma)\le\gap{\Sigma}\le2\norm{\Sigma}_{\op}$.
Both margins are positive exactly when $\Sigma\succ0$: if both are
positive and $x\in\kerop\Sigma$, the supporting inequalities force
$x^\top Jx=0$ and then $(\Sigma-\aplus{\Sigma}J)x=0$, so $x=0$.
The convergence theorem requires only their sum to be positive, so it
also covers some singular covariances. Each positive margin makes the
corresponding principal block positive definite, giving the rank bounds
in Remark~\ref{rem:singular-support}.
\end{remark}

\subsection{Linear Convergence}\label{app:gf}
We prove the spectral-margin theorem and its quantitative trajectory
bounds first. Conditional and determinant-based criteria follow as
supplementary comparisons.

\begin{proposition}[Quantitative trajectory bounds]\label{prop:gf-lift-tail}
Under Theorem~\ref{thm:gf}, set
$a_0:=\lambda_{\min}(\Sigma_0)$, $B_0:=\norm{\Sigma_0}_{\op}$,
$I_0:=\gap{\Sigma_0}^{-1}\sqrt{2\Delta_0/\kappa}$, and $B_\infty:=B_0e^{2I_0}$.
Then
\begin{align}
 a_0e^{-2I_0}\Id_m&\preceq\Sigma_t\preceq B_\infty\Id_m,\label{eq:gf-tube}\\
 \norm{\Sigma_t-\Sigma_\infty}_{\op}
 &\le2B_\infty I_0e^{-\kappa \gap{\Sigma_0}t},\label{eq:gf-tail}\\
 \norm{\Svu_t-\Svu_\infty}_{\F}
 &\le2B_\infty I_0e^{-\kappa \gap{\Sigma_0}t}.\label{eq:gf-predictor-tail}
\end{align}
For the characteristic lift of any initial law in $\cP_2$ with this
second moment, there is a limit $\mu_\infty$ and
\begin{equation}\label{eq:gf-w2}
 W_2(\mu_t,\mu_\infty)
 \le\min\!\left\{\sqrt{\tr\Sigma_0}\,e^{I_0},\sqrt{2B_\infty}\right\}
       I_0e^{-\kappa \gap{\Sigma_0} t}.
\end{equation}
The same bound applies to $\norm{W_t-W_\infty}_{\F}/\sqrt n$ for particles.
The smallest positive covariance eigenvalue also satisfies
$\lambda_{\min}^{+}(\Sigma_t)\ge
e^{-2I_0}\lambda_{\min}^{+}(\Sigma_0)$, including at $t=\infty$.
\end{proposition}

\begin{proof}[Proof of Theorem~\ref{thm:gf} and Proposition~\ref{prop:gf-lift-tail}]
The predictor equation \eqref{eq:block-flow-general} gives the exact
dissipation identity on every positive semidefinite trajectory:
\begin{equation}\label{eq:flow-dissipation}
\frac{\mathrm{d}}{\mathrm{d}t}\ell(\Svu_t)
=-\tr(E_t^\top E_t\Suu_t)-\tr(E_t^\top\Svv_tE_t)\le0.
\end{equation}
The transported hidden-block bounds from Proposition~\ref{lem:margins}
therefore give $-\dot\Delta\ge \gap{\Sigma_0}\norm E_{\F}^2\ge2\kappa \gap{\Sigma_0}\Delta$.
On any finite existence interval, the preceding dissipation and PL give,
where the loss gap is positive,
\[
 -\frac{\dd}{\dd t}\sqrt{\Delta_t}
 \ge\frac{\gap{\Sigma_0}\norm{E_t}_{\F}^2}{2\sqrt{\Delta_t}}
 \ge \gap{\Sigma_0}\sqrt{\kappa/2}\norm{E_t}_{\F}.
\]
At zero loss gap the flow is stationary. Hence, for every finite $T>t$,
\[
 \int_t^T\norm{E_r}_{\F}\,\dd r
 \le\frac1{\gap{\Sigma_0}}\sqrt{\frac2\kappa}
       (\sqrt{\Delta_t}-\sqrt{\Delta_T})
 \le I_0e^{-\kappa \gap{\Sigma_0}t}.
\]
Together with smoothness, this gives
\begin{equation}\label{eq:force-tail}
 \norm{G_t}_{\op}=\norm{E_t}_{\op}
 \le\sqrt{2L\Delta_0}e^{-\kappa \gap{\Sigma_0}t},\qquad
 \int_t^T\norm{G_r}_{\op}\,\dd r\le I_0e^{-\kappa \gap{\Sigma_0}t}.
\end{equation}
Both $\Phi_t$ and $\Phi_t^{-1}$ have norm at most $e^{I_0}$, even when
$\Sigma_0$ is singular. Their congruence representation proves
\eqref{eq:gf-tube} before the maximal existence time. The covariance vector
field is locally Lipschitz on all symmetric matrices, and the upper bound
prevents finite-time blow-up; hence the solution is global.
Integrating $\norm{\dot\Sigma_t}_{\op}\le2B_\infty\norm{G_t}_{\op}$
proves the Cauchy property and \eqref{eq:gf-tail}. Integrability of the
force also makes $\Phi_t$ Cauchy; its limit is invertible since
$\sigma_{\min}(\Phi_t)\ge e^{-I_0}$.
Thus $\Sigma_\infty=\Phi_\infty\Sigma_0\Phi_\infty^\top$ has the same
rank as $\Sigma_0$ and is optimal by \eqref{eq:gf-rate}.
The singular-value bound for congruence gives the claimed positive
eigenvalue floor, without requiring full rank.
Directly in predictor Frobenius norm,
$\norm{\dSvu_t}_{\F}\le2B_\infty\norm{E_t}_{\F}$;
its integral proves \eqref{eq:gf-predictor-tail} without a dimension factor.
Finally $\norm{\Phi_t-\Phi_\infty}_{\op}\le e^{I_0}I_0e^{-\kappa \gap{\Sigma_0}t}$.
Coupling the two characteristic images of the same initial particle gives
the first bound in \eqref{eq:gf-w2}. For the second, use
\[
 \left(\int\norm{G_t w}^2\,\dd\mu_t(w)\right)^{1/2}
 =\sqrt{\tr(G_t\Sigma_tG_t)}
 \le\sqrt{2B_\infty}\,\norm{E_t}_{\F},
\]
and integrate along the same characteristics. The identical estimate
holds for empirical factors after division by $\sqrt n$.
If $\Delta_0=0$, the first-order flow is stationary and
the same formulas hold with $I_0=0$.
\end{proof}
\paragraph{Alternative nondegeneracy criteria}
The same argument can use any independently established hidden-block
lower bounds. The common-bound formulation below uses $2c$ in place of
$\gap{\Sigma_0}$. For $\Sigma_0\succ0$, taking
$c:=\min(\aplus{\Sigma_0},\aminus{\Sigma_0})$ recovers a possibly weaker
rate than Theorem~\ref{thm:gf}.

\begin{lemma}[Conditional hidden-block criterion]
\label{thm:uniform-nondegeneracy-convergence}
Under Assumption~\ref{ass:loss}, start \eqref{eq:bures-G} from
$\Sigma_0\succ0$. If $\Suu_t\succeq c\Id$ and $\Svv_t\succeq c\Id$
for some $c>0$ throughout its maximal existence interval, the flow is
global and
\begin{equation}
\label{eq:uniform-nondegeneracy-energy-rate}
\Delta_t\le\Delta_0e^{-4\kappa ct}.
\end{equation}
Moreover $\Sigma_t$ converges to $\Sigma_\infty\succ0$ with
$\Svu_\infty\in\cS_\star$.
\end{lemma}

Quadratic growth further gives
$\dist(\Svu_t,\cS_\star)^2\le2\Delta_0e^{-4\kappa ct}/\lambda$,
where $\lambda=\kappa$ is always available.

\begin{proof}[Proof of Lemma~\ref{thm:uniform-nondegeneracy-convergence}]
The assumed hidden bounds give
$-\dot\Delta\ge2c\norm E_{\F}^2\ge4c\kappa\Delta$.
On every positive loss gap interval, differentiating $\sqrt{\Delta_t}$ gives
\[
 -\frac{\mathrm{d}}{\mathrm{d}t}\sqrt{\Delta_t}
 \ge c\sqrt{2\kappa}\,\norm{E_t}_{\F},\qquad
 \int_t^T\norm{G_r}_{\op}\,\mathrm{d}r
 \le\frac{\sqrt{\Delta_t}-\sqrt{\Delta_T}}{c\sqrt{2\kappa}}.
\]
The same bound holds through zero loss gap by stationarity and is valid
before the maximal existence time.
Both characteristic matrices $\Phi_t$ and $\Phi_t^{-1}$ are bounded by
the exponential of this integral. Their covariance congruence gives
\begin{equation}\label{eq:uniform-nondegeneracy-covariance-bounds}
 \norm{\Sigma_t}_{\op}\le\norm{\Sigma_0}_{\op}
  \exp\!\left(\frac1c\sqrt{\frac{2\Delta_0}{\kappa}}\right),\qquad
 \lambda_{\min}(\Sigma_t)\ge\lambda_{\min}(\Sigma_0)
  \exp\!\left(-\frac1c\sqrt{\frac{2\Delta_0}{\kappa}}\right).
\end{equation}
These bounds rule out finite-time blow-up. The tail of the integral proves convergence to a
positive-definite covariance with optimal predictor.
Quadratic growth gives the distance bound.
\end{proof}

For coercive losses, determinant conservation and the second spectral
law yield an alternative, generally much coarser, source of such bounds.
\begin{proposition}[Coercivity prevents covariance collapse]
\label{prop:coercive-nondegeneracy}
Under Assumption~\ref{ass:loss}, suppose $\ell$ is coercive and
$\Sigma_0\succ0$. Let $M:=\sup\{\norm S_{\F}:\ell(S)\le\ell(\Svu_0)\}$.
On its maximal existence interval the Bures flow satisfies
$c\Id_m\preceq\Sigma_t\preceq B\Id_m$, where
$B:=\sqrt{\tr((J\Sigma_0)^2)+2M^2}+M$ and
$c:=\det\Sigma_0/B^{m-1}>0$.
In particular both hidden blocks are bounded below by $c\Id$.
\end{proposition}

\begin{proof}[Proof of Proposition~\ref{prop:coercive-nondegeneracy}]
The characteristic congruence \eqref{eq:orbit} preserves positive
definiteness on the maximal existence interval. Loss dissipation gives
$\norm{\Svu_t}_{\F}\le M$. Since $\tr G_t=0$, Jacobi's formula yields
$\det\Sigma_t=\det\Sigma_0>0$. The second spectral law gives
\[
 \norm{\Suu_t}_{\F}^2+\norm{\Svv_t}_{\F}^2
 =\tr((J\Sigma_0)^2)+2\norm{\Svu_t}_{\F}^2
 \le\tr((J\Sigma_0)^2)+2M^2.
\]
Split $\Sigma_t$ into its block-diagonal and off-diagonal parts. Their
operator norms are bounded by the square root of the last expression
and by $\norm{\Svu_t}_{\op}\le M$, respectively. Thus
$\norm{\Sigma_t}_{\op}\le B$. Finally,
$\det\Sigma_0\le\lambda_{\min}(\Sigma_t)B^{m-1}$ gives
$\Sigma_t\succeq c\Id_m$, and the same lower bound holds on both
principal blocks.
\end{proof}

\begin{remark}[Coarse ellipticity versus the spectral-margin rate]
\label{rem:coarse-rate}
At $\Sigma_0=\alpha\Id_m$, the determinant argument gives
\[
c_{\rm coarse}
:=\alpha\left(\frac{\alpha}
{\sqrt{m\alpha^2+2M^2}+M}\right)^{m-1}.
\]
Its rate $4\kappa c_{\rm coarse}$ can be extremely small in dimension,
whereas the support gap $\gap{\Sigma_0}=2\alpha$ gives $4\alpha\kappa$ directly.
For general initial covariances, the margin argument likewise gives a
rate determined by the signed spectrum, with no additional dimension factor.
Trace and initial loss alone cannot provide a uniform positive rate as
the support gap tends to zero
(Example~\ref{ex:slow}).
\end{remark}

\begin{remark}[Degenerate boundary leaves]\label{rem:degenerate-boundary-leaves}
In the $2\times2$ covariance setting of Section~\ref{sec:laws}, write
$\Sigma=\begin{psmallmatrix}p&r\\r&q\end{psmallmatrix}$,
$b_0:=(p_0-q_0)/2$, and $\alpha^2:=p_0q_0-r_0^2$,
with predictor loss $\ell(r):=(r-\beta)^2/2$.
On the boundary \(b_0=\alpha=0\), the scalar equation is
\(\dot r=-2|r|(r-\beta)\).
If \(\beta\ne0\), the ray with \(r_0\beta>0\) converges to \(\beta\),
whereas \(r_0\beta<0\) converges to the nonoptimal apex \(r=0\);
the apex itself is stationary. For example, when \(\beta=1\) and
\(p_0=q_0=1\), \(r_0=-1\), one has
\(p_t=q_t=-r_t=(2e^{2t}-1)^{-1}\).
When \(\beta=0\), one has \(r_t=r_0/(1+2|r_0|t)\): convergence to the
optimal apex is algebraic for \(r_0\ne0\).
These boundary cases have $\gap{\Sigma_0}=0$ and are not covered by the
positive-support-gap convergence theorem.
\end{remark}

\begin{example}[Fixed trace and predictor, arbitrarily slow escape]\label{ex:slow}
Let $d_u=d_v=2$, $0<\delta<1$, and
\[
 \Sigma_0=\diag(\delta,2-\delta,\delta,2-\delta),\qquad
 \ell(S):=\tfrac12\norm{S-e_1e_1^\top}_{\F}^2.
\]
Here $L=\kappa=1$, $\tr\Sigma_0=4$, $\Svu_0=0$, and $\Delta_0=1/2$,
independently of $\delta$. The active scalar block evolves by
\[
 r'=2\sqrt{r^2+\delta^2}(1-r),\qquad r(0)=0.
\]
The instantaneous relative decay of the loss at zero is $4\delta$.
Until $r=1/2$, $r'\le2(r+\delta)$, so
$r(t)\le\delta(e^{2t}-1)$. The time to reach $1/2$ is at least
$\tfrac12\log(1+1/(2\delta))$, which diverges as $\delta\downarrow0$.
There can therefore be no uniform positive exponential rate with a fixed
finite prefactor determined only by the displayed trace, initial loss gap, $L$, and
$\kappa$. At $\delta=0$ the active block is stationary and nonoptimal.
Theorem~\ref{thm:gf} records exactly the missing scale: $\gap{\Sigma_0}=2\delta$.
The slow escape occurs even for a strongly convex loss with a unique
predictor minimizer.
\end{example}


\subsection{Finite Width Convergence}\label{app:sampling}
Write $r_0:=\rankop\tSigma_0$ for the population covariance rank.
\begin{theorem}[Relative initialization stability]\label{thm:relative}
Under Assumption~\ref{ass:loss}, let $\tSigma_0\succeq0$ with
$\gap{\tSigma_0}>0$, and put $\widetilde B_0:=\norm{\tSigma_0}_{\op}$ and
$\widetilde\Delta_0:=\ell(\tSvu_0)-\ell_\star$.
If an empirical covariance satisfies
\begin{equation}\label{eq:relative}
 (1-\varepsilon)\tSigma_0\preceq\Sigma_0
 \preceq(1+\varepsilon)\tSigma_0,\qquad 0<\varepsilon<1,
\end{equation}
then its first-order factor flow is global, converges to an optimal
predictor, and its loss gap $\Delta_t=\ell(\Svu_t)-\ell_\star$ satisfies
\begin{equation}\label{eq:sample-rate}
 \Delta_t\le\Delta_0
 e^{-2(1-\varepsilon)\kappa \gap{\tSigma_0}t}.
\end{equation}
Only the lower inequality in \eqref{eq:relative} is needed for this claim.
With both inequalities, $d_\wedge:=\min(d_u,d_v)$ gives the deterministic
prefactor bound
\begin{equation}\label{eq:sample-prefactor}
 \Delta_0\le
 \left(\sqrt{\widetilde\Delta_0}+\varepsilon \widetilde B_0\sqrt{Ld_\wedge/2}\right)^2.
\end{equation}
For the sampling law in Corollary~\ref{cor:finite-width}, the width
condition \eqref{eq:width} guarantees \eqref{eq:relative} with probability
at least $1-\delta$.
\end{theorem}


\begin{proof}[Proof of Theorem~\ref{thm:relative}]
From $\tSigma_0-\aplus{\tSigma_0}J\succeq0$ and the lower relative inequality,
$\Sigma_0-(1-\varepsilon)\aplus{\tSigma_0}J\succeq0$, and similarly for
$\aminus{\tSigma_0}$. Transporting these inequalities along the empirical
fundamental matrix allows us to apply the proof of Theorem~\ref{thm:gf} with margin
sum $(1-\varepsilon)\gap{\tSigma_0}$ and initial empirical loss gap. The two-sided relative
comparison gives
$\norm{\Sigma_0-\tSigma_0}_{\op}\le\varepsilon \widetilde B_0$ and hence
$\norm{\Svu_0-\tSvu_0}_{\F}\le\sqrt{d_\wedge}\varepsilon \widetilde B_0$.
Equation~\eqref{eq:sample-prefactor} follows from \eqref{eq:perturb}.
For the probability bound, choose an orthonormal basis
$T\in\R^{m\times r_0}$ of the range of $\tSigma_0$ and write
$\tSigma_0=T\Lambda T^\top$, where $\Lambda\succ0$.
The vectors $\eta_i:=T^\top\xi_i\in\R^{r_0}$ have identity second
moment and subgaussian norm at most $K$, and
$w_i=T\Lambda^{1/2}\eta_i$.
The covariance tail bound~\citep[Exercise 4.7.3]{Vershynin2018}
applied in this range gives
\[
 \norm{n^{-1}\textstyle\sum_i\eta_i\eta_i^\top-\Id_{r_0}}_{\op}
 \le CK^2\left(\sqrt{\frac{r_0+u}{n}}+\frac{r_0+u}{n}\right)
\]
with failure probability at most $2e^{-u}$. Set $u:=\log(2/\delta)$ and
adjust $C$ in \eqref{eq:width}. Multiplication on the left by
$T\Lambda^{1/2}$ and on the right by its transpose gives
\eqref{eq:relative}.
\end{proof}
\begin{proof}[Proof of Corollary~\ref{cor:finite-width}]
Theorem~\ref{thm:relative} gives the relative covariance event and the
loss rate. Its lower bound gives
$\gap{\Sigma_0}\ge(1-\varepsilon)\gap{\tSigma_0}>0$,
so Theorem~\ref{thm:gf} gives global existence and factor convergence.
The two-sided event preserves the kernel of $\tSigma_0$, so the empirical
initial covariance is positive definite exactly when $\tSigma_0$ is.
\end{proof}

\paragraph{Small width under full-rank isotropic initialization}
The rank improvement in Remark~\ref{rem:rank-adaptive-width} concerns
the chosen population law. A positive population support gap alone does not
guarantee a positive empirical support gap at small width when
$\tSigma_0=\sigma^2\Id_m$, $\sigma>0$.
\begin{example}[Isotropic Gaussian sampling can lose the support gap]
\label{ex:white-small-width}
Let $d_u=D$, $d_v=1$, and sample from $N(0,\tSigma_0)$ with
$\tSigma_0=\sigma^2\Id_{D+1}$, $\sigma>0$.
There is a universal $c>0$ such that, if
$1\le n\le D/16$, then
$\mathbb P\{\gap{\Sigma_0}=0\}\ge1-4e^{-cD}$,
although $\gap{\tSigma_0}=2\sigma^2$.
Thus a width bound depending only on $\min(d_u,d_v)$ cannot uniformly
preserve a positive support gap when sampling from $N(0,\sigma^2\Id_m)$. Convergence of
these trajectories is not decided by this example.
\end{example}
\begin{proof}
Write $U_0=\sigma X$ and $V_0=\sigma y^\top$, where
$X\in\R^{D\times n}$ and $y\in\R^n$ have independent standard
Gaussian entries. Gaussian singular-value and norm concentration
\citep[Section~7.3]{Vershynin2018}, with deviation $\sqrt D/8$, gives
\[
 \sigma_{\min}(X)\ge\sqrt D-\sqrt n-\sqrt D/8\ge5\sqrt D/8,
 \qquad \norm y\le\sqrt n+\sqrt D/8\le3\sqrt D/8
\]
with probability at least $1-4e^{-cD}$, for a universal $c>0$.
Hence $X^\top X-yy^\top\succeq(D/4)\Id_n$.
The nonzero spectrum of $J\Sigma_0$ is the spectrum of
$\sigma^2(X^\top X-yy^\top)/n$, so it has $n$ positive eigenvalues
and $D+1-n\ge2$ zero eigenvalues. In \eqref{eq:gap-spectrum},
$d_v=1$ therefore gives $\gap{\Sigma_0}=\lambda_2-\lambda_1=0$.
\end{proof}

\paragraph{A separate initial-accuracy budget}
A relative covariance bound preserves the rate; an absolute bound also
controls the initial loss. In the sampling setting of
Corollary~\ref{cor:finite-width}, write
$\widetilde B_0:=\norm{\tSigma_0}_{\op}$, $\widetilde\Delta_0:=\ell(\tSvu_0)-\ell_\star$, and
$d_\wedge:=\min(d_u,d_v)$.
Besides $\varepsilon,\delta\in(0,1)$, choose an absolute tolerance
$\zeta>0$ and set
\begin{equation}\label{eq:sampling-budgets}
 \beta:=\min\{\varepsilon,\zeta/\widetilde B_0\},\qquad
 \overline\Delta:=
 \left(\sqrt{\widetilde\Delta_0}+\zeta\sqrt{Ld_\wedge/2}\right)^2.
\end{equation}

\begin{theorem}[High-probability finite-width convergence]
\label{thm:sampling-budgets}
Under Assumption~\ref{ass:loss}, let $\tSigma_0\succeq0$ with
$\gap{\tSigma_0}>0$ and use the sampling law and tolerances above.
There is a universal $C>0$ such that, if
\begin{equation}\label{eq:budget-width}
 n\ge CK^4\beta^{-2}\bigl(r_0+\log(2/\delta)\bigr),
\end{equation}
then with probability at least $1-\delta$ the empirical covariance has
positive support gap, the factor flow is global, and
\[
 \Delta_t\le\Delta_0e^{-2\kappa(1-\varepsilon)\gap{\tSigma_0}t}
 \le\overline\Delta e^{-2\kappa(1-\varepsilon)\gap{\tSigma_0}t}.
\]
The factors converge, their limiting predictor belongs to $\cS_\star$, and
$\dist(\Svu_t,\cS_\star)^2\le2\Delta_t/\kappa$.
\end{theorem}

\begin{proof}[Proof of Theorem~\ref{thm:sampling-budgets}]
Apply the relative concentration event with tolerance $\beta$.
It implies the relative bounds with tolerance $\varepsilon$ and
$\norm{\Sigma_0-\tSigma_0}_{\op}\le\beta \widetilde B_0\le\zeta$.
The two initial supporting inequalities give the margin bounds, and
Lemma~\ref{lem:pl} gives the loss gap bound. More precisely, putting
$\underline g:=(1-\varepsilon)\gap{\tSigma_0}$ and $\overline B_0:=(1+\varepsilon)\widetilde B_0$,
the same event gives
\begin{equation}\label{eq:sampling-event}
 a_\pm(\Sigma_0)\ge(1-\varepsilon)a_\pm(\tSigma_0),\quad
 \norm{\Sigma_0}_{\op}\le\overline B_0,\quad
 \Delta_0\le\overline\Delta.
\end{equation}
Theorem~\ref{thm:gf} then applies
to the empirical covariance. The distance estimate follows from
Lemma~\ref{lem:pl}. More generally, any valid quadratic-growth constant
$\lambda$ gives $\dist(\Svu_t,\cS_\star)^2\le2\Delta_t/\lambda$.
\end{proof}
\begin{proposition}[The Gaussian initial-error cost]\label{prop:gaussian}
For independent $u_i\sim N(0,\sigma^2\Id_{d_u})$ and
$v_i\sim N(0,\sigma^2\Id_{d_v})$, and the quadratic loss
$\ell(S):=\norm{S-\Svu_\star}_{\F}^2/2$, one has exactly
\begin{equation}\label{eq:gaussian-error}
 \mathbb E\Delta_0
 =\tfrac12\norm{\Svu_\star}_{\F}^2+\frac{\sigma^4d_ud_v}{2n}.
\end{equation}
\end{proposition}

For balanced dimensions and fixed $\sigma$, width $n$ proportional to $m$
therefore gives an expected initial sampling error of order $m$, despite
the dimension-free exponent $4(1-\varepsilon)\sigma^2\kappa$ at any
fixed relative tolerance $\varepsilon\in(0,1)$. Keeping that error bounded
in expectation requires $n$ of order $d_ud_v$, when the teacher norm is
fixed. Thus width proportional to $m$ can preserve the convergence rate
while allowing the expected initial error to grow with dimension.

\begin{proof}[Proof of Proposition~\ref{prop:gaussian}]
Each entry of $n^{-1}\sum_i v_i u_i^\top$ has mean zero and variance
$\sigma^4/n$. Sum these variances; the cross term with the teacher has
zero expectation.
\end{proof}

\subsection{Discrete-Time Convergence}\label{app:gd}
The factor update transports the initial supporting inequalities even
though it no longer preserves $J$ exactly. To track this error, define
\begin{equation}\label{eq:gd-signature}
\Phi_0:=\Id,\quad\Phi_{k+1}:=M_k\Phi_k,\quad
\mathcal J_k:=\Phi_kJ\Phi_k^\top,\quad\mathcal D_k:=\mathcal J_k-J.
\end{equation}
The cross block of \eqref{eq:gd-covariance} also gives
\begin{equation}\label{eq:gd-predictor-update}
 \Svu_{k+1}-\Svu_k
 =-\tau(E_k\Suu_k+\Svv_kE_k)+\tau^2E_k(\Svu_k)^\top E_k.
\end{equation}
The next two lemmas control the accumulated defect $\mathcal D_k$ and
the quadratic correction in the predictor update.

\begin{lemma}[Transported support and accumulated defect]\label{lem:gd-support}
If $\Sigma_0\succeq0$, with initial margins $\aplus{\Sigma_0},\aminus{\Sigma_0}$ as in
\eqref{eq:margins}, then every finite iterate satisfies
\[
 \Sigma_k-\aplus{\Sigma_0}\mathcal J_k\succeq0,\qquad
 \Sigma_k+\aminus{\Sigma_0}\mathcal J_k\succeq0.
\]
If $\norm{\mathcal D_k}_{\op}\le\eta$ for some $\eta\in[0,1)$, then
$\Suu_k\succeq(1-\eta)\aplus{\Sigma_0}\Id$ and
$\Svv_k\succeq(1-\eta)\aminus{\Sigma_0}\Id$. With $e_k:=\norm{E_k}_{\F}$,
\begin{equation}\label{eq:gd-defect-recursion}
 \mathcal D_{k+1}=M_k\mathcal D_kM_k+\tau^2G_kJG_k,\qquad
 \norm{\mathcal D_N}_{\op}
 \le\tau\exp\left(2\sum_{j<N}\tau e_j\right)
             \sum_{j<N}\tau e_j^2.
\end{equation}
\end{lemma}

\begin{proof}[Proof of Lemma~\ref{lem:gd-support}]
The proof of Proposition~\ref{lem:margins} gives the initial supporting
inequalities. Congruence by $\Phi_k$ transports them, regardless of
invertibility. The $uu$ and $vv$ blocks of $\mathcal J_k$ are respectively
at least $(1-\eta)\Id$ and at most $-(1-\eta)\Id$, proving the hidden
lower bounds. Since $G_kJ+JG_k=0$, one has
$M_kJM_k=J+\tau^2G_kJG_k$, giving the recursion. Its norm is at most
$(1+\tau e_k)^2\norm{\mathcal D_k}_{\op}+\tau^2e_k^2$.
Unroll from zero and bound the amplification products by
$\exp(2\sum_{j<N}\tau e_j)$.
\end{proof}
\paragraph{Descent and length.}
The bootstrap needs two estimates: loss descent and a finite total
gradient length. The next lemma obtains both from the same step condition.
For the proof, abbreviate the descent scale, initial gradient bound, and
covariance bound as
\begin{equation}\label{eq:gd-constants}
 \alpha_\star:=\frac{\gap{\Sigma_0}}{2\sqrt2},\qquad
 M_{\nabla}:=\sqrt{2L\Delta_0},\qquad
 B_\star:=\norm{\Sigma_0}_{\op}e^{2I_\star}.
\end{equation}
Then $I_\star=\sqrt{2\Delta_0/\kappa}/\alpha_\star$ and
$q=1-2\alpha_\star\kappa\tau$. Condition~\eqref{eq:gd-step} is
equivalent to requiring both
$\tau\le\alpha_\star/(B_\star M_{\nabla}+\frac92LB_\star^2)$ and
$\tau\le(1-1/\sqrt2)\alpha_\star/(e^{2I_\star}\Delta_0)$.
\begin{lemma}[Descent with unequal hidden lower bounds]\label{lem:gd-descent}
Suppose $\norm{\Sigma_k}_{\op}\le B$, $\Delta_k\le\Delta_0$,
$\Suu_k\succeq b_+\Id$, and $\Svv_k\succeq b_-\Id$.
Write $h:=b_++b_->c>0$. If $\tau M_{\nabla}\le1$ and
$\tau(BM_{\nabla}+\frac92LB^2)\le h-c$, then
\begin{equation}\label{eq:gd-descent}
 \Delta_{k+1}\le\Delta_k-c\tau e_k^2
 \le(1-2c\kappa\tau)\Delta_k.
\end{equation}
If this descent holds for $j=r,\ldots,N-1$, then
\begin{equation}\label{eq:gd-length}
 \sum_{j=r}^{N-1}\tau e_j^2\le\Delta_r/c,\qquad
 \sum_{j=r}^{N-1}\tau e_j
 \le\frac1c\sqrt{\frac2\kappa}
            (\sqrt{\Delta_r}-\sqrt{\Delta_N}).
\end{equation}
\end{lemma}

\begin{proof}[Proof of Lemma~\ref{lem:gd-descent}]
Let $Z:=\Svu_{k+1}-\Svu_k$. By \eqref{eq:gd-covariance} and smoothness,
$\norm Z_{\F}\le3\tau Be_k$ and
$\ip{E_k}{Z}\le-\tau h e_k^2+\tau^2BM_{\nabla}e_k^2$.
The descent lemma therefore gives the first assertion. Telescope it for
the squared-gradient sum. At positive loss gap, PL gives
\[
 \sqrt{\Delta_j}-\sqrt{\Delta_{j+1}}
 \ge\frac{c\tau e_j^2}{2\sqrt{\Delta_j}}
 \ge c\tau\sqrt{\kappa/2}\,e_j.
\]
At zero loss gap $e_j=0$. Summation proves the length bound.
\end{proof}

The proof below also establishes the following explicit tails in
Theorem~\ref{thm:gd}:
\begin{equation}\label{eq:gd-tails}
 \frac{\norm{W_k-W_\infty}_{\F}}{\sqrt n}
 \le\sqrt{2B_\star}\,I_\star q^{k/2},\qquad
 \norm{\Sigma_k-\Sigma_\infty}_{\op}
 \le2B_\star I_\star q^{k/2}.
\end{equation}

\begin{proof}[Proof of Theorem~\ref{thm:gd}]
Let $\eta:=1-1/\sqrt2$. Since $\gap{\Sigma_0}\le2B_0$ for
$B_0:=\norm{\Sigma_0}_{\op}$, one has
$\alpha_\star\le B_0/\sqrt2\le B_\star/\sqrt2$.
The first restriction in \eqref{eq:gd-step} implies
\[
 \tau M_{\nabla}<1/\sqrt2,\qquad
 0<2\alpha_\star\kappa\tau
 \le\frac{4\kappa\alpha_\star^2}{9LB_\star^2}\le\frac29.
\]
In particular $q\in[7/9,1)$.
Inductively suppose, through iterate $N$, that
$\Delta_k\le q^k\Delta_0$,
$\norm{\Sigma_k}_{\op}\le B_\star$, and
$\norm{\mathcal D_k}_{\op}\le\eta$, with descent proved for $k<N$.
These assertions start at $N=0$ since $\mathcal D_0=0$.
At $N$, Lemma~\ref{lem:gd-support} supplies hidden lower bounds with sum
$(1-\eta)\gap{\Sigma_0}=2\alpha_\star$. Lemma~\ref{lem:gd-descent} applies with
$c=\alpha_\star$, $h=2\alpha_\star$, and $B=B_\star$, proving the next
loss-gap estimate. Its finite-horizon sums now give
\[
 \sum_{j=0}^{N}\tau e_j\le I_\star,\qquad
 \sum_{j=0}^{N}\tau e_j^2\le\Delta_0/\alpha_\star.
\]
Thus $\norm{\Phi_{N+1}}_{\op}\le e^{I_\star}$ implies the next covariance
upper bound, and \eqref{eq:gd-defect-recursion} gives
$\norm{\mathcal D_{N+1}}_{\op}
\le\tau e^{2I_\star}\Delta_0/\alpha_\star\le\eta$.
This closes the induction.

Taking tails in \eqref{eq:gd-length} yields
$\sum_{j\ge k}\tau e_j\le I_\star q^{k/2}$.
Factor increments have norm
\[
 \norm{W_{j+1}-W_j}_{\F}/\sqrt n
=\tau\sqrt{\tr(G_j\Sigma_jG_j)}
\le\tau\sqrt{2B_\star}\,e_j.
\]
This proves factor convergence and its tail estimate. Similarly
$\norm{\Phi_k-\Phi_\infty}_{\op}\le e^{I_\star}I_\star q^{k/2}$
gives the covariance tail bound by subtracting the two congruences.
The limit predictor is optimal by continuity and the loss-gap bound.
Finally, if $r:=\tau M_{\nabla}<1$, every factor update is invertible and
\[
 \sigma_{\min}(\Phi_k)\ge\exp\left(-\frac{I_\star}{1-r}\right),\qquad
 \lambda_{\min}^{+}(\Sigma_k)\ge\lambda_{\min}^{+}(\Sigma_0)
 \exp\left(-\frac{2I_\star}{1-r}\right).
\]
Indeed $\sigma_{\min}(M_j)\ge1-\tau e_j>0$ and
$\log(1-x)\ge-x/(1-r)$ for $0\le x\le r$. The bound makes
$\Phi_\infty$ invertible, proving rank preservation also at the limit.
\end{proof}
The factor-increment estimate also applies to population discrete
pushforward updates. Coupling each particle with its limiting position
gives a Wasserstein tail at most $\sqrt{2B_\star}I_\star q^{k/2}$.

\paragraph{Isotropic initialization and local stability}
For an exact isotropic empirical second moment
$\Sigma_0=W_0W_0^\top/n=\sigma^2\Id_m$ and $\Delta_0=\ell(0)-\ell_\star>0$,
$I_\star=2\sqrt{\Delta_0}/(\sigma^2\sqrt\kappa)$ and the explicit
condition is
\[
 \tau\le\min\left\{
 \frac1{2e^{2I_\star}\sqrt{L\Delta_0}
                 +\frac9{\sqrt2}L\sigma^2e^{4I_\star}},\quad
 \frac{(\sqrt2-1)\sigma^2}{2e^{2I_\star}\Delta_0}
 \right\}.
\]
The contraction factor is $1-\sqrt2\kappa\sigma^2\tau$.
Here the empirical covariance is exactly $\sigma^2\Id_m$. For iid Gaussian
initialization, its fluctuations are handled by Corollary~\ref{cor:sampled-gd}.
As $\Delta_0\downarrow0$ with $\sigma,L,\kappa$ fixed, the sufficient
threshold tends to $\sqrt2/(9L\sigma^2)$.


\begin{remark}[Local stability and the global stepsize]\label{rem:gd-local}
Suppose $\ell$ is $C^2$ near a minimizer at zero, and linearize at
$\Sigma=\sigma^2\Id$. The predictor perturbation evolves to first order by
$\delta\Svu\mapsto(\Id-2\tau\sigma^2\nabla^2\ell(0))[\delta\Svu]$.
Thus $0<\tau<1/(L\sigma^2)$ contracts its nonzero-curvature directions;
kernel directions are neutral at this order. Theorem~\ref{thm:gd} additionally controls
covariance growth and the accumulated signature defect along the entire
trajectory, explaining its more restrictive global threshold.
\end{remark}

\paragraph{Adjustable rate budgets.}
Smaller steps permit a smaller loss in the conserved conditioning,
bringing the discrete exponent closer to its continuous-time value.
\begin{proposition}[Adjustable rate budgets]\label{prop:gd-adjustable}
Under the loss and initialization assumptions of Theorem~\ref{thm:gd},
choose $0<c<h<\gap{\Sigma_0}$ and set
$I_c:=\sqrt{2\Delta_0/\kappa}/c$, $B_c:=B_0e^{2I_c}$.
The convergence and rank conclusions of Theorem~\ref{thm:gd} hold with
contraction factor $q_c:=1-2c\kappa\tau$, and the tail bounds
\eqref{eq:gd-tails} hold with $(I_\star,B_\star,q)$ replaced by
$(I_c,B_c,q_c)$, if
\[
 0<\tau\le\min\left\{\frac1{2M_{\nabla}},\
 \frac{h-c}{B_cM_{\nabla}+\frac92LB_c^2},\
 \frac{(1-h/\gap{\Sigma_0})c}{e^{2I_c}\Delta_0}\right\}.
\]
\end{proposition}

Taking $c=(1-\epsilon)\gap{\Sigma_0}$, $h=(1-\epsilon/2)\gap{\Sigma_0}$ for any
$\epsilon\in(0,1)$ certifies exponent $2\kappa(1-\epsilon)\gap{\Sigma_0}$ in time
$k\tau$, at an explicit sufficiently small step. This is arbitrarily
close to the flow exponent.

\begin{proof}[Proof of Proposition~\ref{prop:gd-adjustable}]
Use $\eta=1-h/\gap{\Sigma_0}$ in the preceding induction and
Lemma~\ref{lem:gd-descent} with $c,h$. The positivity guard follows from
$c(h-c)\le \gap{\Sigma_0}^2/4\le B_0^2$: $2c\kappa\tau\le4/9$.
The remaining estimates use $I_c,B_c$.
\end{proof}
\begin{corollary}[A stepsize fixed before sampling]\label{cor:sampled-gd}
Use the population parameters and width condition of
Theorem~\ref{thm:sampling-budgets}, with the bounds in
\eqref{eq:sampling-event}, and define before sampling
\[
 \overline\alpha_\star:=\frac{\underline g}{2\sqrt2},\quad
 \overline M_{\nabla}:=\sqrt{2L\overline\Delta},\quad
 \overline I_\star:=\frac1{\overline\alpha_\star}
                    \sqrt{\frac{2\overline\Delta}{\kappa}},\quad
 \overline B_\star:=\overline B_0e^{2\overline I_\star}.
\]
Choose a deterministic positive stepsize satisfying
\begin{equation}\label{eq:sampled-gd-step}
 \tau\le\min\left\{
 \frac{\overline\alpha_\star}
 {\overline B_\star\overline M_{\nabla}+\frac92L\overline B_\star^2},\quad
 \frac{(1-1/\sqrt2)\overline\alpha_\star}
 {e^{2\overline I_\star}\overline\Delta}\right\}.
\end{equation}
With probability at least $1-\delta$, factor GD converges to an optimal
predictor, and
\[
 \Delta_k
 \le(1-2\overline\alpha_\star\kappa\tau)^k\Delta_0
 \le(1-2\overline\alpha_\star\kappa\tau)^k\overline\Delta.
\]
Its factor/covariance tail estimates hold with the corresponding barred
constants and contraction factor.
\end{corollary}

At fixed relative tolerance, the sampled stepsize can deteriorate with
the dimension through $\overline\Delta$. The independent absolute
accuracy budget controls this initial-loss contribution.


\begin{proof}[Proof of Corollary~\ref{cor:sampled-gd}]
On \eqref{eq:sampling-event}, the empirical descent scale is at least
$\overline\alpha_\star$, whereas the empirical loss gap, force bound, length
bound, and covariance upper bound are at most their barred counterparts.
Both admissible stepsize terms increase with the descent scale and
decrease with the loss gap and initial covariance norm. Hence the chosen
step satisfies Theorem~\ref{thm:gd} for the empirical initialization.
Its contraction factor is no larger than the stated one. If the realized
initial loss gap vanishes, the empirical factors are stationary, and the
conclusions hold directly.
\end{proof}

\paragraph{Boundary cases.}
For the scalar loss $\ell(s):=(s+1)^2/2$, the singular initialization
$U_0=V_0\ne0$ preserves equality of the factors and keeps
$\Svu_k=\norm{U_k}^2/n\ge0$, so its loss gap is at least $1/2$.
Conversely, even full-rank isotropic initialization does not justify arbitrary steps:
for $\ell(s):=(s-1)^2/2$, $\Sigma_0=\Id_2$, and $\tau=1$, one gets
$\Sigma_1=\bigl(\begin{smallmatrix}2&2\\2&2\end{smallmatrix}\bigr)$
and $\Sigma_2=0$, after which training is trapped.

\subsection{Deep Linear Residual Networks}\label{app:resnets}
Write
$P_{b:a}:=A_b\cdots A_a$ when $a\le b$, and $P_{b:a}:=\Id$ otherwise.
Thus $P_{>s}=P_{S-1:s+1}$ and $P_{<s}=P_{s-1:0}$.
Differentiating $P_S=P_{>s}A_sP_{<s}$ gives
\begin{equation}
  \label{eq:resnet-layer-gradient}
\nabla_{\Svu_s}\ell(P_S)
=h\widetilde E_s,
\qquad
\widetilde E_s=P_{>s}^\top E P_{<s}^\top.
\end{equation}
In residual mean-field time the particle equations are
$(\dot u_{s,i},\dot v_{s,i})
=-(\widetilde E_s^\top v_{s,i},\widetilde E_su_{s,i})$.
The cross block of \eqref{eq:resnet-layer-bures} is therefore
$\dSvu_s=-\widetilde E_s\Suu_s-\Svv_s\widetilde E_s$.

In the finite-depth setting of Theorem~\ref{thm:resnet-finite}, put
$\overline{\sigma}_S:=(1+B/S)^{S-1}$ and
$I_{\rm res}:=\overline{\sigma}_S(\overline g_S\underline{\sigma}_S^2)^{-1}\sqrt{2\Delta_0/\kappa}$.
The proof gives the common layer-force tail
\begin{equation}\label{eq:resnet-length}
 \int_t^\infty\norm{\widetilde E_s(r)}_{\F}\,\dd r
 \le I_{\rm res}e^{-\kappa\overline g_S\underline{\sigma}_S^2t}.
\end{equation}

\begin{proof}[Proof of Theorem~\ref{thm:resnet-finite}]
Write $g_s:=\gap{\Sigma_s(0)}$ within this proof.
The endpoint derivative and layer covariance equations give
\[
 -\dot\Delta
 =\frac1S\sum_{s=0}^{S-1}
 \left[\tr(\widetilde E_s^\top \widetilde E_s\Suu_s)+\tr(\widetilde E_s\widetilde E_s^\top\Svv_s)\right]
 \ge\frac1S\sum_{s=0}^{S-1}g_s\norm{\widetilde E_s}_{\F}^2.
\]
Both the lower and upper singular-value bounds of the residual factors
multiply through the prefix and suffix, so
$\underline{\sigma}_S\norm E_{\F}\le\norm{\widetilde E_s}_{\F}\le \overline{\sigma}_S\norm E_{\F}$.
Therefore
\[
 -\dot\Delta\ge\overline g_S\underline{\sigma}_S^2\norm E_{\F}^2
 \ge2\kappa\overline g_S\underline{\sigma}_S^2\Delta.
\]
This proves the loss bound on each finite existence interval using the
average layer margin. The square-root PL
argument used in Theorem~\ref{thm:gf}, followed by the upper force bound,
gives \eqref{eq:resnet-length} first with any finite endpoint.
Thus all layer fundamental matrices and their inverses have norm at most
$e^{I_{\rm res}}$. Their covariances remain bounded and positive definite.
The finite-dimensional vector field is locally Lipschitz; these bounds
prevent a finite maximal existence time. The force-tail bound now holds
on $[t,\infty)$ and implies layer convergence by the argument of
Theorem~\ref{thm:gf}. Continuity of the finite residual product proves
optimality of its limit.
\end{proof}

\begin{proof}[Proof of Theorem~\ref{thm:resnet-continuous}]
The $L^1$ training equation in Section~\ref{sec:resnets} means
\begin{equation}\label{eq:resnet-integral}
 \Sigma_s(t)=\Sigma_s(0)-\int_0^t
 (\widetilde G_s(r)\Sigma_s(r)+\Sigma_s(r)\widetilde G_s(r))\,\dd r,
 \qquad\widetilde G_s=G_{\widetilde E_s}.
\end{equation}
This equality in $L^1$ gives an absolutely continuous representative for
almost every layer. The path bound makes the forces uniformly bounded
over depth on finite training intervals, justifying differentiation
of the endpoint and integration over depth.
An $L^1$ residual field has invertible propagators. Gronwall applied to
each propagator and its inverse gives
$e^{-B}\norm E_{\F}\le\norm{\widetilde E_s}_{\F}\le e^B\norm E_{\F}$.
The depth integrals of the conserved supporting inequalities yield
\[
 -\dot\Delta
 =\int_0^1\left[\tr(\widetilde E_s^\top \widetilde E_s\Suu_s)
                    +\tr(\widetilde E_s\widetilde E_s^\top\Svv_s)\right]\,\dd s
 \ge\overline g e^{-2B}\norm E_{\F}^2.
\]
These integrals are finite by the $L^1$ covariance assumption and the
uniform-in-depth force bounds. PL proves \eqref{eq:continuous-resnet-rate}.
As $\gap{\Sigma_s(0)}>0$ almost everywhere and
$\gap{\Sigma_s(0)}\le2\norm{\Sigma_s(0)}_{\op}$, its integral is finite and positive.
The square-root argument gives the common layer-force tail bound
\[
 \int_t^\infty\norm{\widetilde E_s(r)}_{\F}\,\dd r
 \le I_{\rm res}e^{-\kappa\overline g e^{-2B}t},\qquad
 I_{\rm res}:=\frac{e^{3B}}{\overline g}
                         \sqrt{\frac{2\Delta_0}{\kappa}}.
\]
The fundamental-matrix argument proves almost-everywhere convergence
and positive definiteness. The covariance tails are bounded by
$2e^{2I_{\rm res}}\norm{\Sigma_s(0)}_{\op}I_{\rm res}
e^{-\kappa\overline g e^{-2B}t}$, an integrable bound over depth.
This proves $L^1$ convergence. Continuity of the endpoint map on bounded
$L^1$ sets, obtained by variation of constants and Gronwall, proves
optimality of the limiting endpoint.
\end{proof}

\begin{remark}[Sharpness of the boundedness hypothesis]
\label{rem:resnet-path-escape}
Unlike the two-layer flow of Theorem~\ref{thm:gf}, positive conserved
layer margins do not rule out unbounded weight growth as $S\to\infty$
or along continuous-depth training. A linear neural ODE has positive
determinant, so it cannot fit an orientation-reversing target; training
may instead approach a singular endpoint through unbounded residuals.
For a concrete example, take the invertible scalar target $p_\star:=-1$ and loss
$\ell(p):=(p+1)^2/2$, with $L=\kappa=1$. Initialize every layer with
$\Sigma_s(0)=\Id_2$, realizable at width $n=2$; every conserved support
gap equals $2$. At each finite depth $S\ge2$, all residuals decrease
to $-S$, so $\sup_{t\ge0}\max_s|\Svu_s(t)|=S$ and no uniform $B<S$
exists. In continuous depth they decrease without bound, giving
$\sup_{t\ge0}\int_0^1|\Svu_s(t)|\,\dd s=\infty$.
Both flows are global, but their endpoints tend to zero and their losses
to $1/2>\ell_\star=0$. At finite depth the target is representable:
failure occurs along this symmetric trajectory, whereas the neural ODE
excludes the target altogether. In both cases, positive conserved margins
therefore fail to ensure the residual bounds required by the theorems.
\end{remark}

\begin{proof}[Verification of Remark~\ref{rem:resnet-path-escape}]
The scalar target $-1$ is invertible, and
$|\ell'(p)|^2=(p+1)^2=2\ell(p)$, so $L=\kappa=1$ and
$\ell_\star=0$. The width-two initialization
$U_s(0)=(\sqrt2,0)$, $V_s(0)=(0,\sqrt2)$ gives
$\Sigma_s(0)=\Id_2$ and initial endpoint $1$ at every depth.
For scalar layers write $a_s:=\Suu_s=\Svv_s$ and
$r_s:=\Svu_s$. Equality of the diagonal blocks is preserved, and the
layer equations reduce to
\[
 \dot a_s=-2\widetilde E_s r_s,\qquad
 \dot r_s=-2\widetilde E_s a_s.
\]
Thus $a_s^2-r_s^2=1$, and positivity gives
\[
 \Sigma_s(t)=\begin{pmatrix}
  \sqrt{1+r_s(t)^2}&r_s(t)\\
  r_s(t)&\sqrt{1+r_s(t)^2}
 \end{pmatrix}.
\]
The eigenvalues of $J\Sigma_s$ are $\pm1$, whereas those of
$\Sigma_s$ are $\sqrt{1+r_s^2}\pm|r_s|>0$.
Uniqueness and symmetry preserve identical layers. Denote their common
residual by $r_S$. At finite depth,
$\widetilde E_s=(p_S+1)(1+r_S/S)^{S-1}$, giving
\begin{equation}\label{eq:resnet-escape-finite}
 \begin{aligned}
 p_S&=(1+r_S/S)^S,\qquad r_S(0)=0,\\
 \dot r_S&=-2\sqrt{1+r_S^2}\big((1+r_S/S)^S+1\big)(1+r_S/S)^{S-1}.
 \end{aligned}
\end{equation}
For $S\ge2$, its vector field is locally Lipschitz and vanishes at $r_S=-S$.
Starting at zero, the solution decreases and cannot reach this equilibrium
in finite time. It stays in the compact interval $[-S,0]$, so it is global.
Any limit strictly above $-S$ would have strictly negative velocity;
hence its limit is $-S$, $p_S\to0$, and $\ell(p_S)\to1/2$.
In particular, every candidate bound $B<S$ is exceeded at a finite
training time, although $|r_S(t)|<S$ at every finite time.

The finite-depth architecture can represent $-1$ by taking one residual
factor equal to $-1$ and the others equal to $1$. The obstruction here
is the preserved symmetry of the initialized trajectory: all its
residual factors stay positive and approach zero. At $S=1$, this
obstruction disappears; $\dot r_1=-2\sqrt{1+r_1^2}(r_1+2)$ instead
gives $r_1\to-2$ and the optimal endpoint $-1$.

For a depth-constant field, the continuous propagators are
$P_{b,a}=e^{(b-a)r}$, so $\widetilde E_s=(e^r+1)e^r$ and
\begin{equation}\label{eq:resnet-escape-continuous}
 p=e^r,\qquad \dot r=-2\sqrt{1+r^2}(e^r+1)e^r,\qquad r(0)=0.
\end{equation}
For $r\le0$ its velocity is bounded, strictly negative, and has no finite
zero. The solution is therefore global and tends to $-\infty$.
Since $\int_0^1|\Svu_s(t)|\,\dd s=|r(t)|$, every finite integral
bound is eventually exceeded as well. The endpoint tends to zero.
The limiting loss $1/2$ is the infimum over positive scalar endpoints
and exceeds the ambient minimum $\ell_\star=0$.
The associated depth-constant covariance is a global locally absolutely
continuous $L^1$ trajectory, positive definite at every finite time.
Moreover, for $S\ge2$ and $-S<r_S\le0$,
$(1+r_S/S)^{S-1}\le e^{r_S/2}$ and $(1+r_S/S)^S+1\le2$.
Since $\sqrt{1+r^2}e^{r/2}\le1$ for $r\le0$, both $|\dot r_S|$ and
$|\dot r|$ are at most $4$. On $[0,T]$, all residuals therefore lie in
$[-4T,0]$. The finite-depth vector fields converge uniformly on this
compact interval, and continuous dependence gives $r_S\to r$ uniformly
on every bounded training interval $[0,T]$.
\end{proof}

\begin{remark}[Zero-target variant]
The same path bounds can fail while the loss tends to zero. With
$\ell(p):=p^2/2$ and the same initialization, the equations are
$\dot r_S=-2\sqrt{1+r_S^2}(1+r_S/S)^{2S-1}$ and
$\dot r=-2\sqrt{1+r^2}e^{2r}$. The preceding argument gives
$r_S\to-S$ for every $S\ge1$, $r\to-\infty$ in continuous depth,
and endpoints tending to the zero target. A uniform residual-path
bound would keep the endpoints bounded away from zero instead.
\end{remark}

\begin{remark}[Residual-path control and orientation]
Positive-definite layer covariances can coexist with singular residual
products. For a scalar depth-two network, take both layer predictors equal to $-2$,
so $A_0=A_1=0$, and the endpoint loss $\ell(p):=(p-1)^2/2$.
Both backpropagated gradients vanish although the endpoint is not optimal.
These predictors admit positive-definite layer covariances
$\bigl(\begin{smallmatrix}\sqrt5&-2\\-2&\sqrt5\end{smallmatrix}\bigr)$,
whose signed eigenvalues are $+1,-1$. Extending the earlier discretization and sampling arguments to this setting
requires control of the residual products along the corresponding training
trajectory; the two-layer GD theorem does not apply verbatim.

In continuous depth, $\det P(t)=
\exp(\int_0^1\tr\Svu_s(t)\,\dd s)>0$.
The uniform path bound also keeps $\sigma_{\min}(P(t))\ge e^{-B}$,
so any convergent endpoint remains in the orientation-preserving
invertible component. If every global minimizer has negative determinant,
these path bounds are incompatible with convergence to a global minimizer.
\end{remark}

\subsection{Heavy-ball Momentum}\label{app:momentum-proof}
We first expand the closed phase equation, then establish global
existence, position-deformation estimates, and the energy bootstrap.
For positions $x:=(u,v)$ and velocities $y:=(p,q)$, let
$\nu_t\in\cP_2(\R^{2m})$ be their joint distribution. Its covariance is
$\Phase_t:=\int\binom xy\binom xy^{\!\top}\,\dd\nu_t(x,y)$,
with mixed block $\Xi_t:=\int xy^\top\,\dd\nu_t(x,y)$.
Finite-width factors correspond to
$\nu_t:=n^{-1}\sum_i\delta_{(x_i(t),y_i(t))}$.
The characteristics satisfy
\begin{equation}
\label{eq:single-particle-momentum-main}
\dot x=y,\qquad
\dot y=-G(\Svu)x-\gamma y.
\end{equation}
Taking its second moments gives the block system
\begin{equation}
\label{eq:phase-block-ode}
\begin{aligned}
\dSpos_t&=\Xi_t+\Xi_t^\top,\qquad
\dot\Xi_t=\Svel_t-\Spos_tG_t-\gamma\Xi_t,\\
\dSvel_t&=-G_t\Xi_t-\Xi_t^\top G_t-2\gamma\Svel_t.
\end{aligned}
\end{equation}
In particular, the predictor velocity and acceleration are
\begin{equation}\label{eq:predictor-speed}
 \dSvu=\Sigma^{qu}+\Sigma^{vp},\qquad
 \ddSvu+\gamma\dSvu=2\Sqp-E\Suu-\Svv E.
\end{equation}

\begin{proof}[Proof of Proposition~\ref{prop:phase-laws}]
Since \(G_t\) is symmetric,
\(\Mob_t^\top \Jcan+\Jcan\Mob_t=-\gamma \Jcan\).
Differentiate the two-characteristic pairing and the fundamental-matrix
identity. Invertibility also gives
\(\Psi_t\Jcan(\Psi_t)^\top=e^{-\gamma t}\Jcan\), whence
\(\Jcan\Phase_t=e^{-\gamma t}
(\Psi_t)^{-\top}(\Jcan\Phase_0)(\Psi_t)^\top\).
\end{proof}

\begin{proposition}[Global full phase flow]\label{prop:global-phase}
If $\ell$ is $C^1$, has locally Lipschitz gradient, and is bounded below,
every $\Phase_0\succeq0$ has a global positive semidefinite solution of
\eqref{eq:phase-block-ode}.  This assertion alone gives neither uniform position ellipticity
nor convergence to a minimizer.
\end{proposition}

\begin{proof}[Proof of Proposition~\ref{prop:global-phase}]
Local well-posedness follows from local Lipschitz continuity. The linear
phase characteristic lift preserves positivity. The energy identity gives
$\tr\Svel_t\le2\cE_0$. Positivity and Cauchy--Schwarz yield
\[
 \sqrt{\tr\Spos_t}
 \le\sqrt{\tr\Spos_0}+t\sqrt{2\cE_0},\qquad
 \norm{\Xi_t}_{\F}^2\le\tr\Spos_t\tr\Svel_t.
\]
The first inequality also follows directly from the $L^2$ triangle
inequality along characteristics, including times when the trace is zero.
Every block is bounded on each finite interval; continuation proves
global existence.
\end{proof}
For the matrix calculations in the momentum sections, abbreviate
$A:=\Spos$ and $B:=\Svel$. Whenever $A\succ0$, define
\begin{equation}\label{eq:regression}
 D:=\Xi^\top A^{-1},\qquad Q:=B-DAD^\top\succeq0.
\end{equation}
Here $D$ is a regression matrix, not assumed symmetric, and $Q$ is its
positive semidefinite residual covariance (the Schur complement).
They encode all phase blocks by $\Xi=AD^\top$ and $B=DAD^\top+Q$.
\begin{proposition}[General covariance coordinates]\label{prop:regression}
On the chart $A\succ0$, \eqref{eq:phase-block-ode} is equivalent to
\begin{equation}\label{eq:regression-ode}
 \begin{aligned}
 \dot A&=DA+AD^\top,\\
 \dot D&=-G-\gamma D-D^2+QA^{-1},\\
 \dot Q&=-2\gamma Q-DQ-QD^\top.
 \end{aligned}
\end{equation}
Moreover, positive semidefiniteness controls the position deformation:
\begin{equation}\label{eq:regression-bound-main}
 \norm D_{\op}^2\le\frac{2\mathcal{K}}{\lambda_{\min}(A)}.
\end{equation}
\end{proposition}


\begin{proof}[Proof of Proposition~\ref{prop:regression}]
Differentiate $D=\Xi^\top A^{-1}$ and $Q=B-DAD^\top$, using
\eqref{eq:phase-block-ode}. This gives the three displayed identities.
Since $DAD^\top\preceq B$, its trace is at most $2\mathcal{K}$, whereas
$\lambda_{\min}(A)\norm D_{\op}^2\le\tr(DAD^\top)$.
\end{proof}
\begin{lemma}[Conditional phase-energy decay]\label{lem:conditional}
Under Assumption~\ref{ass:loss}, let a positive semidefinite solution of
\eqref{eq:phase-block-ode} satisfy $a\Id_m\preceq\Spos_t\preceq b\Id_m$ for
$0\le t\le T$, where $0<a\le b<\infty$. Put
$\mathfrak K:=4a\kappa$, $\mathfrak L:=4bL$, and
$\rho:=\mathfrak L+2\sqrt{2L\cE_0}$. Then
\begin{equation}\label{eq:conditional}
 \cE_t\le3\cE_0 e^{-\lambda t},\qquad
 \lambda:=\frac{a\kappa\gamma}{\rho+2\gamma^2},\qquad 0\le t\le T.
\end{equation}
\end{lemma}
\begin{proof}
Set $\cN:=\tr(G\Spos G)$, the mean squared particle force.
The hidden principal blocks are between $a\Id$ and $b\Id$, hence
\[
 2a\norm E_{\F}^2\le\cN\le2b\norm E_{\F}^2.
\]
Combining PL with the smooth-gradient upper bound in
Lemma~\ref{lem:pl}, and applying Cauchy--Schwarz to the phase law, gives
\begin{equation}\label{eq:conditional-force-bounds}
 \mathfrak K\Delta\le \cN\le\mathfrak L\Delta,\qquad
 |\dot\Delta|=|\tr(G\Xi)|\le\sqrt{2\mathcal{K} \cN}.
\end{equation}
All second derivatives below are understood almost everywhere. Lipschitz
continuity of $\nabla\ell$ implies that $E_t$ is locally absolutely
continuous and $\norm{\dot E_t}_{\F}\le L\norm{\dSvu_t}_{\F}$.
Differentiating \eqref{eq:predictor-speed} gives
\[
 \ddot\Delta=H-\cN-\gamma\dot\Delta,\qquad
 H:=\ip{\dot E}{\dSvu}+2\ip E{\Sqp}\le\rho\mathcal{K}.
\]
For the last bound, positive phase blocks give
$\norm{\Sigma^{qu}}_{\F}^2\le b\tr\Sigma^q$ and
$\norm{\Sigma^{vp}}_{\F}^2\le b\tr\Sigma^p$, hence
$\norm{\dSvu}_{\F}^2\le4b\mathcal{K}$. Also
$\norm{\Sqp}_{\F}\le\sqrt{\tr\Sigma^q\tr\Sigma^p}\le\mathcal{K}$ and
$\norm E_{\F}\le\sqrt{2L\cE_0}$.

Choose $\epsilon:=\gamma/[2(\rho+2\gamma^2)]$ and
$\mathcal V:=\cE+\epsilon\dot\Delta$. Since
\[
 \epsilon\le\frac1{4\sqrt{2\rho}}
 \le\frac1{\sqrt{2\mathfrak L}},\qquad
 |\dot\Delta|\le\sqrt{\mathfrak L/2}\,\cE,
\]
one has $\cE/2\le\mathcal V\le3\cE/2$.
The inequality $\gamma\sqrt{2\mathcal{K} \cN}\le \cN/4+2\gamma^2\mathcal{K}$ yields
\[
 \dot{\mathcal V}\le-\tfrac34\epsilon \cN
       -\bigl(2\gamma-\epsilon(\rho+2\gamma^2)\bigr)\mathcal{K}
 \le-\tfrac34\epsilon\mathfrak K\Delta-\gamma\mathcal{K}.
\]
Now $\tfrac34\epsilon\mathfrak K\le\gamma$, since $a\kappa\le bL$
and $\rho\ge4bL$. It follows that
$\dot{\mathcal V}\le-\tfrac23(\tfrac34\epsilon\mathfrak K)\mathcal V
=-\lambda\mathcal V$.
Integrating the a.e. inequality proves \eqref{eq:conditional}.

\end{proof}

Under Theorem~\ref{thm:momentum}, the position-covariance tail is
\begin{equation}\label{eq:momentum-tail}
 \norm{\Spos_t-\Spos_\infty}_{\F}
 \le\frac{4\sqrt{6b_s\cE_0}}{\lambda_s(\gamma)}
       e^{-\lambda_s(\gamma)t/2}.
\end{equation}
For the position factors $W$ or the position marginal of the characteristic
phase law, the normalized Frobenius or $W_2$ tail has prefactor
$2\sqrt{6\cE_0}/\lambda_s(\gamma)$ at the same exponent.

\begin{proof}[Proof of Theorem~\ref{thm:momentum}]
The full phase flow is global by Proposition~\ref{prop:global-phase}.
On the connected position-positive chart containing time zero, let
$h(t):=\int_0^t\norm{D_r}_{\op}\,\dd r$.
The equation $\dot A=DA+AD^\top$ yields
$A_t=Z_t A_0Z_t^\top$ with $\dot Z=DZ$, $Z_0=\Id_m$, and thus
\[
 a_0e^{-2h(t)}\Id_m\preceq A_t\preceq b_0e^{2h(t)}\Id_m.
\]
As long as $h\le s$, Lemma~\ref{lem:conditional} and
Proposition~\ref{prop:regression} imply
\[
 h(t)\le\int_0^t\sqrt{2\mathcal{K}_r/a_s}\,\dd r
 \le\frac{2}{\lambda_s}\sqrt{\frac{6\cE_0}{a_s}}
 =\mathfrak S_s(\gamma)<s.
\]
There is no first time at which $h=s$. Nor can the chart terminate earlier:
on each such interval $A\succeq a_s\Id$ and the globally defined full
phase blocks stay bounded on finite horizons. This proves the global tube
and energy estimate. Furthermore the positive-block factorization gives
$\norm\Xi_{\F}^2\le b_s\tr B\le6b_s\cE_0 e^{-\lambda_st}$.
Integrate $\dot A=\Xi+\Xi^\top$ to get \eqref{eq:momentum-tail}.
For particles or a characteristic phase law, the remaining $L^2$ position displacement is
at most $2\sqrt{6\cE_0}\,e^{-\lambda_st/2}/\lambda_s$.
The limit is therefore attained, with zero loss gap and full phase
covariance $\diag(A_\infty,0)$.
\end{proof}

If $\cE_0=0$, the nonnegative kinetic and potential energies both vanish,
the force vanishes, and the momentum trajectory is stationary. Positive
definiteness of the \emph{phase} covariance is never required; in particular
zero-speed initialization is covered.

\subsection{An exact decomposition of hidden kinetic energy}
Prescribing the position covariance and its derivative does not determine
the kinetic energy: velocity components that leave the covariance
unchanged can still contribute to it.

\begin{proposition}[Minimum-energy lift of a covariance velocity]\label{prop:kinetic}
Fix $A\succ0$ and a symmetric matrix $H$ as the desired value of $\dot A$.
Let $C=C^\top$ be the unique solution of $CA+AC=H$.
Every positive semidefinite phase covariance with this position moment and
derivative has regression coordinates
\[
 D=C+\Omega,\qquad \Omega A+A\Omega^\top=0,\qquad Q\succeq0,
\]
and its kinetic energy decomposes as
\begin{equation}\label{eq:kinetic-decomposition}
 \mathcal{K}=\tfrac12\tr(CAC)
    +\tfrac12\tr(\Omega A\Omega^\top)+\tfrac12\tr Q.
\end{equation}
Thus the graph lift $D=C$, $Q=0$ has minimum kinetic energy.
An additional positive semidefinite Schur complement contributes exactly
$\tr Q/2$.
\end{proposition}
\begin{proof}
In an eigenbasis of $A$, the Sylvester equation divides $H_{ij}$ by the
strictly positive sum of eigenvalues, proving existence and uniqueness.
Equation $DA+AD^\top=H$ is equivalent to the stated condition on $\Omega$.
Since $B=DAD^\top+Q$, expand its trace. The two cross terms vanish in sum:
$\Omega A$ is skew-symmetric and $C$ is symmetric. Each term in
\eqref{eq:kinetic-decomposition} is nonnegative.
\end{proof}

\begin{example}[Identical position covariance and covariance derivative, arbitrarily large energy]
Let $A=\sigma^2\Id_m$, $C=0$, choose any skew-symmetric $\Omega$, and set
$D=\Omega$, $Q=0$. The initial predictor is zero because
$A=\sigma^2\Id_m$, and $\dot A=0$ makes its derivative zero as well,
independently of $\Omega$. Nevertheless
$\mathcal{K}=\sigma^2\norm\Omega_{\F}^2/2$ can be arbitrarily large. The phase
covariance is positive semidefinite; nonzero $\Omega$ gives a nonsymmetric regression matrix.
For fixed $a_0,b_0,L,\kappa$, the best feasibility ratio
\eqref{eq:budget} tends to zero as $\cE_0\to\infty$. Thus the initial
position covariance and predictor speed do not determine whether the
energy-damping certificate holds.
\end{example}

%%% Local Variables:
%%% mode: LaTeX
%%% TeX-master: "../main"
%%% End:
